\section{The energy geometry at fixed mass and momentum}
\label{sec:r12-constraint}

The spectral reduction for a specified smooth branch requires control of
the real linearized kernel. The variational assertion below does not:
the imaginary response alone forces a negative constrained energy
curvature when the homogeneous order exceeds two. Every component of
momentum is preserved, and no differentiability of the dilation group on
the individual profile is assumed.

Let $s=m/2$, $\alpha=2m/N$, and let $a$ be a positive, real-even,
elliptic homogeneous symbol of degree $m>1$. Write
\[
 E(u)=\frac12 A(u)-\frac1{\alpha+2}\|u\|_{\alpha+2}^{\alpha+2},
 \quad M(u)=\|u\|_2^2,
 \quad P(u)=\operatorname{Im}\int\overline u\nabla u.
\]
Let $Q$ be a real sharp ground profile at frequency one. The scalar
phase-coercivity theorem gives
\[
 L_-=a(D)+1-|Q|^\alpha\geq0,\qquad \ker L_-=\mathbb RQ,
\]
with coercivity on $H^s\cap Q^\perp$. Define
\begin{equation}\label{eq:r12-phase-T}
 Z_j=L_-^{-1}\partial_jQ,\qquad
 T_{ij}=\langle\partial_iQ,Z_j\rangle.
\end{equation}
The inverse and pairing can be understood in energy-space form sense.
In particular $T$ is positive definite.

\begin{theorem}[Constrained curvature without a profile branch]
\label{thm:r12-constrained-curvature}
Suppose $v_n\to0$, $v_n\ne0$, and exact profiles satisfy
\[
 (a(D)+1+iv_n\cdot\nabla)U_n-|U_n|^\alpha U_n=0,
 \qquad U_n\longrightarrow Q\quad\hbox{strongly in }H^s,
\]
after phase and translation. For every sufficiently large $n$ there is
a continuous $H^s$ curve $V_n(\lambda)$, defined near $\lambda=1$, with
\[
 V_n(1)=U_n,\qquad
 M(V_n(\lambda))=M(U_n),\qquad
 P(V_n(\lambda))=P(U_n).
\]
Its scalar energy is $C^2$ and obeys
\begin{equation}\label{eq:r12-constrained-curvature}
 \left.\frac{d}{d\lambda}E(V_n(\lambda))\right|_1=0,
 \qquad
 \left.\frac{d^2}{d\lambda^2}E(V_n(\lambda))\right|_1
 =(2-m)v_n^TTv_n+o(|v_n|^2).
\end{equation}
For $m>2$ the profile is therefore not a local energy minimum at its
fixed mass and momentum: both sides of this curve have strictly lower
energy. The statement assumes neither a differentiable family of
profiles nor nondegeneracy of the real block $L_+$.
\end{theorem}

\begin{proof}
We suppress the sequence index and keep errors relative to its limit.
Choose the phase so that $U=X+iY$ and $\langle Y,Q\rangle=0$.
The imaginary profile equation is
\[
 L_-Y=-v\cdot\nabla X+(|U|^\alpha-|Q|^\alpha)Y.
\]
Put $p=\alpha+2$. Since $H^s\hookrightarrow L^p$, the potential
$|U|^\alpha-|Q|^\alpha$ tends to zero as a quadratic form on $H^s$:
this follows from H\"older and its convergence in $L^{p/\alpha}$.
The derivative maps $H^s$ continuously to $H^{-s}$ because $2s=m>1$.
Project onto $Q^\perp$ and apply phase coercivity. First
$\|Y\|_{H^s}=O(|v|)$; substituting this bound back into the equation gives
\begin{equation}\label{eq:r12-imaginary-response}
 Y=-\sum_jv_jZ_j+o(|v|)\quad\hbox{in }H^s.
\end{equation}
No estimate for a derivative of the real profile is needed.
The bilinear momentum form is continuous on $H^s$, and
$P_j(U)=-2\langle Y,\partial_jX\rangle$. Hence
\begin{equation}\label{eq:r12-P-response}
 P(U)=2Tv+o(|v|).
\end{equation}
The positivity of $T$ follows from phase coercivity and the fact that a
nonzero $L^2$ function cannot have a vanishing nonzero directional
derivative.

For $c=(c_0,c_1,\ldots,c_N)\in\mathbb R^{N+1}$ set
\[
 W(c)=(1+c_0)U+i\sum_{j=1}^N c_jZ_j.
\]
The Jacobian at $c=0$ of $c\mapsto(M(W(c)),P(W(c)))$ is
\[
 \begin{pmatrix}
  2M(U)&O(|v|)\\
  2P(U)&-2T+o(1)
 \end{pmatrix}.
\]
It is invertible for sufficiently small velocity. The finite-dimensional
implicit-function theorem supplies a smooth $c(\lambda)$ with $c(1)=0$
and
\begin{equation}\label{eq:r12-W-charges}
 M(W(\lambda))=M(U),\qquad
 P(W(\lambda))=\lambda^{-1}P(U).
\end{equation}
Here and below $W(\lambda)$ denotes $W(c(\lambda))$.
Differentiating these equations and using
\eqref{eq:r12-P-response} gives
\[
 c_j'(1)=v_j+o(|v|),\qquad c_0'(1)=O(|v|^2),
 \qquad h:=W'(1)=i\sum_jv_jZ_j+o(|v|)\quad\hbox{in }H^s.
\]
The amplitude estimate uses
$\langle U,iZ_j\rangle=\langle Y,Z_j\rangle=O(|v|)$.

Define
\[
 V(\lambda,x)=\lambda^{N/2}W(\lambda,\lambda x).
\]
Critical homogeneity gives the exact identities
\begin{equation}\label{eq:r12-scaled-charges}
 M(V(\lambda))=M(U),\quad P(V(\lambda))=P(U),\quad
 E(V(\lambda))=\lambda^mE(W(\lambda)).
\end{equation}
Strong continuity of dilation gives a continuous $H^s$ curve. Although
that curve need not be differentiable in $H^s$, its scalar energy is
$C^2$ by the last identity. This is the reason that no weighted moment
or dilation-domain assumption enters the theorem.

Let $\mathcal L_U$ be the real Hessian of
$S=E+M/2-v\cdot P/2$, and put
$q_U(h)=\langle h,\mathcal L_Uh\rangle$ in form notation.
Because $S'(U)=0$, differentiating \eqref{eq:r12-W-charges} yields
\[
 \left.\frac{d}{d\lambda}E(W(\lambda))\right|_1=-\frac12v\cdot P(U),
 \qquad
 \left.\frac{d^2}{d\lambda^2}E(W(\lambda))\right|_1
 =q_U(h)+v\cdot P(U).
\]
Using the exact profile identity $v\cdot P=2mE$ in
\eqref{eq:r12-scaled-charges}, we conclude that
\begin{equation}\label{eq:r12-exact-curvature}
 E(V)'(1)=0,\qquad
 E(V)''(1)=q_U(h)-\frac{m-1}{2}v\cdot P(U).
\end{equation}
The energy is a $C^2$ functional on $H^s$, since $p>2$ and
$H^s\hookrightarrow L^p$. Its Hessian forms therefore converge on bounded
sets as $U\to Q$, $v\to0$. The imaginary block at $Q$ is $L_-$, so
\[
 q_U(h)=q_Q\Big(i\sum_jv_jZ_j\Big)+o(|v|^2)
 =v^TTv+o(|v|^2).
\]
Together with \eqref{eq:r12-P-response} and
\eqref{eq:r12-exact-curvature}, this proves
\eqref{eq:r12-constrained-curvature}. The sign for $m>2$ follows from
positive definiteness of $T$.
\end{proof}

\begin{corollary}[All sufficiently slow action ground waves]
\label{cor:r12-all-ground-waves}
At frequency one, define
\[
 K_v(u)=A(u)+M(u)-v\cdot P(u),\qquad
 \kappa_v=\inf_{u\ne0}\frac{K_v(u)}{\|u\|_p^2},
 \qquad p=2+\frac{2m}{N}.
\]
For sufficiently small $|v|$, this infimum is attained. Normalize each
minimizer $U$ by $K_v(U)=\|U\|_p^p$; these are precisely the action
minimizers on the nonzero Nehari manifold. If $m>2$, then every such
ground wave at every sufficiently small nonzero velocity is not a local
energy minimum at fixed mass and momentum. This conclusion requires no
uniqueness, differentiable branch, or real-profile nondegeneracy.
\end{corollary}

\begin{proof}
Ellipticity and H\"older in Fourier variables give
\[
 |P(u)|\leq C A(u)^{1/m}M(u)^{1-1/m}
          \leq C'(A(u)+M(u)).
\]
Consequently $K_v$ is uniformly equivalent to the $H^s$ norm squared
for sufficiently small $|v|$. The exponent $p>2$ is strictly below the
energy-space Sobolev endpoint, including when $s\geq N/2$.
For completeness, normalize a minimizing sequence by $\|u_n\|_p=1$.
It is bounded in $H^s$ and cannot vanish locally, by the standard
localized Sobolev interpolation argument. After a translation it has a
nonzero weak limit $u$. Local compact embedding and the
Brezis--Lieb identity split the $p$th power norm into $c$ and $1-c$,
where $c=\|u\|_p^p>0$. Weak Hilbert-space orthogonality splits $K_v$.
Thus
\[
 \kappa_v\geq\kappa_v\big(c^{2/p}+(1-c)^{2/p}\big).
\]
Since $2/p<1$, this forces $c=1$. Norm convergence and the coercivity of
$K_v$ give strong $H^s$ convergence and attainment. Rescaling a minimizer
to the Nehari manifold gives the profile equation and
\begin{equation}\label{eq:r12-ground-norm}
 K_v(U)=\|U\|_p^p=\kappa_v^{p/(p-2)}.
\end{equation}
On that manifold $S_v(U)=(1/2-1/p)\|U\|_p^p$, proving the asserted
minimizing equivalence.

The Fourier bound also gives $\kappa_v\to\kappa_0$ as $v\to0$.
Let $v_n\to0$ and let $U_n$ be any Nehari-normalized minimizers.
They are bounded in $H^s$ by \eqref{eq:r12-ground-norm}, and
$K_0(U_n)-K_{v_n}(U_n)=v_n\cdot P(U_n)\to0$. Therefore their normalized
versions form a minimizing sequence for $\kappa_0$.
The preceding compactness proof gives, after translations and passing to
a subsequence, $U_n\to Q$ strongly in $H^s$, where $Q$ is a
frequency-one ground profile.

We verify that this limiting profile is one of the sharp critical
optimizers used in Theorem~\ref{thm:r12-constrained-curvature}.
Write $r=m/N$ and let the sharp critical inequality be
\[
 \|u\|_p^p\leq C_* M(u)^r A(u).
\]
For fixed $A+M$, the product $M^rA$ is maximal at $M=rA$.
A critical dilation of a sharp optimizer attains that ratio, so it also
attains $\kappa_0$. Conversely, attainment of $\kappa_0$ forces equality
in both inequalities, hence sharp critical optimization and $M=rA$.
The density phase-rigidity theorem makes this optimizer real after a
constant phase. Nehari normalization then gives
$\|Q\|_p^p=(1+r)A(Q)$, and equality in
$\|Q\|_p^p\leq(1+r)(M(Q)/M_c)^r A(Q)$ gives $M(Q)=M_c$.

If the asserted uniform small-velocity conclusion failed, there would
be a sequence of nonzero $v_n\to0$ and ground waves lacking a negative
fixed-charge curve. The compactness just proved produces a real limiting
$Q$, to which Theorem~\ref{thm:r12-constrained-curvature} applies. Its
strictly negative curvature contradicts that failure. No uniform
choice of a differentiable family was used.
\end{proof}

\paragraph{Why the phase correction matters.}
The linear boost
\[
 V_b(\lambda,x)=e^{i(1-\lambda)k\cdot x}
                  \lambda^{N/2}U(\lambda x),\qquad
 k=\frac{P(U)}{M(U)},
\]
also fixes both charges. However its curvature, when $a$ has the
homogeneous $C^2$ derivative bounds and $m>2$, is
\begin{equation}\label{eq:r12-linear-boost}
 E(V_b)''(1)=\frac12\left(k^TH_Uk-(m-1)v\cdot P(U)\right),
 \qquad H_U=\int D^2a(\xi)|\widehat U(\xi)|^2\,d\xi.
\end{equation}
Indeed the kinetic term is
$\int a(k+\lambda(\xi-k))|\widehat U(\xi)|^2\,d\xi$.
Euler's identities and the profile boost identity
$\int\nabla a|\widehat U|^2=vM$ give
\eqref{eq:r12-linear-boost}. Near $Q$ its leading term is
\[
 \frac{2}{M(Q)^2}v^TTH_QTv-(m-1)v^TTv.
\]
Phase positivity alone gives no sign for this expression.
By contrast the optimal imaginary correction minimizes
$\langle\psi,L_-\psi\rangle$ under
$-2\langle\partial_jQ,\psi\rangle=-P_j$. Its minimum is
$P^TT^{-1}P/4$, attained at
$\psi=\frac12\sum_j(T^{-1}P)_jZ_j$.
This is the correction used to leading order in the theorem.

\begin{remark}[Energetic and dynamical conclusions]
The theorem gives energy-decreasing perturbations at precisely the same
conserved charges as the wave. It does not alone prove that their
solutions escape a neighborhood of the symmetry orbit. In particular,
Morse index one does not remove an uncontrolled extra kernel: a negative
direction can belong to a zero Jordan chain. Nonlinear or spectral
instability requires the additional dynamical argument stated in its
own theorem. For $1<m<2$ the displayed curve instead has positive
leading curvature, which also does not establish stability.
\end{remark}
